% frontmatter/lupa/verano.tex -- instrucciones para el adulto del cuaderno
% de verano de "Leo con lupa" (ediciones/lupa-verano.tex, ver
% notes/07-ediciones.md): van en vez de frontmatter/lupa/como-usar.tex,
% que habla del año entero. Detrás, "Cómo lee un detective" (el mismo
% del libro entero) y, en vez del mapa del curso, el mapa del verano
% (\mapaEdicion, preamble.tex).
\section*{Cómo usar este cuaderno de verano}

Este cuaderno es el verano de \textit{Leo con lupa}, el tercer cuaderno
de \textit{Aprendo a leer}: sus últimas 13 semanas, en un cuaderno
aparte, para seguir leyendo un poco cada día de las vacaciones. Está
pensado para una niña o un niño que ya lee con soltura y ahora aprende a
leer con atención: textos de unas 120 palabras, con oraciones largas,
de varias partes unidas por \textit{que}, \textit{porque},
\textit{aunque}, \textit{cuando}, \textit{si} o \textit{para que}, y
cada día algo que descubrir en el texto -- no lo que parece que dice,
sino lo que dice de verdad.

Tiene una página para cada día de lunes a viernes, 65 en total, y cada
página sigue siempre el mismo orden:

\begin{itemize}[leftmargin=6mm, itemsep=1mm]
\item \textbf{Fecha}, \textbf{Leído} y \textbf{Notas}, para el adulto.
\item El \textbf{texto del día}, en el recuadro verde, en párrafos. El
  título del recuadro dice cómo leerlo.
\item Una \textbf{actividad de análisis}: para hacerla hay que volver al
  texto. La regla de oro es que \textbf{la respuesta siempre está en el
  texto} -- el de ese día o, si lo dice la instrucción, el de los días
  anteriores de la misma semana. Si no la encuentra, no se la digas:
  pregúntale «¿dónde lo dice?» y deja que vuelva a leer.
\end{itemize}

\textbf{Las actividades.} \textbf{Dibuja con lupa}: dibujar exactamente
lo que describe el texto; no importa que el dibujo sea bonito, sino que
sea exacto (cuántos, de qué color, dónde), y al terminar se repasa con
las preguntas de abajo, texto en mano. \textbf{Mapa}: seguir unas
indicaciones en una cuadrícula y marcar un sitio o un camino.
\textbf{Tabla de pistas}: un problema de lógica; se pone ✓ o ✗ en cada
casilla hasta que solo queda una posibilidad. \textbf{Caza los
errores}: un resumen del texto con fallos que hay que encontrar y
corregir. \textbf{Verdadero o falso}: y, si es falso, cómo es de verdad.
\textbf{Mensaje secreto}: descifrar un mensaje con un código que explica
el texto. \textbf{Responde}, \textbf{Ordena}, \textbf{Relaciona},
\textbf{Ficha}, \textbf{Compara}, \textbf{Adivina}, \textbf{Viñetas} y
\textbf{Crea}: preguntas de «¿por qué?» y «¿cómo lo sabes?», sucesos que
ordenar, datos que sacar de un texto informativo, dos cosas que
comparar, adivinanzas, una historia en tres dibujos y un poco de
escritura propia.

\textbf{Un caso cada semana.} Las páginas cuentan el verano de
\textbf{Lucía}, su hermano \textbf{Dani} y su perro \textbf{Toby} en el
pueblo de la abuela Rosa. Lucía tiene un club de detectives, el
\textbf{Club de la Lupa}, con sus amigos Sofía y Hugo; pero ellos pasan
el verano fuera, así que en el pueblo el club son Lucía, Dani y Toby
-- y Martín, el amigo que Dani tiene en el pueblo. Casi todas las
semanas hay un misterio: las pistas aparecen de lunes a jueves, y el
viernes toca resolverlo --
\textbf{Resuelve el caso}: rodear la respuesta y escribir qué pistas lo
demuestran. La solución se confirma el lunes siguiente, y está en la
\textbf{clave de respuestas} del final del cuaderno. Mejor no mirarla
antes de que lo intente. Y el misterio grande del verano es un mapa del
tesoro que la abuela Rosa dibujó cuando tenía ocho años.

\textbf{El ritmo.} Diez o quince minutos al día. Si el texto se le hace
largo, leedlo juntos en voz alta primero, un párrafo cada uno, y que
luego lo relea sola. Lo que importa es que siga siendo un juego: los
detectives no tienen prisa, tienen cuidado. En la página siguiente está
cómo lee un detective, y después, el verano semana a semana, con un sol
para colorear cada día que se lee; al final, un diploma.
\newpage
